\section{Conclusion}\label{sec:conclusion}

This paper presented \scl{}, a control layer that needs no training, adjusts
itself from a detector's scores, and adapts a frozen tracker through a
declared host contract. The layer splits the logits of recent frames with
nested Otsu thresholds and decides whether the stream is clean or noisy. On
clean streams it keeps the tracker's settings; on noisy streams it chooses
the candidates, assigns them new scores, and sets the thresholds. A
continuation rescue gives single-stage trackers the step they lack.

\scl{} was developed from always-on banding (\aob{}), an earlier design that
acted in every frame, and each step was compared with that baseline. One
frozen configuration was then tested on seven published trackers, four
detection sources, and two datasets with development and predeclared
external roles. It raised the HOTA of OC-SORT and of the predeclared
external PD-SORT by 0.61 points, restored trackers whose thresholds no longer
matched a shifted score scale, raised the MOTA of two-stage trackers with a
noisy detector by 6.0 and 7.9 points, and kept well-matched two-stage
pipelines at their own accuracy. It costs 2.95 ms per frame on a CPU. These
results support placing a self-calibrating layer between the detector and
the tracker in systems where the detector changes while the tracker stays
the same.
